\section{Supplement on extensions of algebras}
\label{section:0.18}

\oldpage[0\textsubscript{IV-1}]{51}

This paragraph assembles a certain number of functorial constructions on rings, which will be used repeatedly in \textsection\textsection~19 and 20; it contains no non-trivial result.

\subsection{Inverse images of augmented rings}
\label{subsection:0.18.1}

\begin{env}[18.1.1]
\label{0.18.1.1}
Given a ring $A$ (not necessarily commutative), the category of \emph{$A$-rings} has as objects the pairs $(B,\rho)$ formed of a ring $B$ and a homomorphism of rings $\rho:A\to B$, for morphisms (also called \emph{$A$-homomorphisms}) the homomorphisms of rings $u:B\to C$ such that, if $\rho:A\to B$ and $\sigma:A\to C$ are the homomorphisms (called \emph{structural}) defining the structure of $A$-ring on $B$ and $C$ respectively, the diagram
\begin{equation}
\tag{18.1.1.1}
  \label{0.18.1.1.1}
  \xymatrix{
    B\ar[r]^u & C\\
    & A\ar[lu]^\rho\ar[u]_\sigma
  }
\end{equation}
is commutative.
The \emph{kernel} $\mathfrak{J}$ of $u$ is a bilateral ideal that is an $A$-\emph{bimodule} (for $\rho$).

If $f:A'\to A$ is a homomorphism of rings, for every $A$-ring $(B,\rho)$, $(B,\rho\circ f)$ is an $A'$-ring, and if $u:B\to C$ is an $A$-homomorphism from the $A$-ring $(B,\rho)$ to the $A$-ring $(C,\sigma)$, it is also an $A'$-homomorphism from the $A'$-ring $(B,\rho\circ f)$ to the $A'$-ring $(C,\sigma\circ f)$; one thus defines a functor from the category of $A$-rings into that of $A'$-rings.
\end{env}

\begin{env}[18.1.2]
\label{0.18.1.2}
Let $B$, $E$, $F$ be three $A$-rings, $f:E\to B$, $g:F\to B$ two $A$-homomorphisms.
Recall that one calls \emph{fibre product} of $E$ and $F$ over $B$ (for the homomorphisms $f$ and $g$) and one denotes $E\times_B F$ the subring of the product $E\times F$ formed of the pairs $(x,y)$ such that $f(x)=g(y)$; the restrictions $p_1:G\to E$, $p_2:G\to F$ of the projections $\operatorname{pr}_1$ and $\operatorname{pr}_2$ to $G$ are still called the \emph{canonical projections};
the structure of $A$-ring of $G$ is defined by the homomorphism $\alpha\mapsto(\rho(\alpha),\sigma(\alpha))$ (where $\rho:A\to E$ and $\sigma:A\to F$ are the structural homomorphisms), which indeed maps $A$ into $G$ by virtue of \sref{0.18.1.1}.
The characteristic property of $E\times_B F$ is that, for every pair of $A$-homomorphisms $u:C\to E$, $v:C\to F$ such that $f\circ u=g\circ v$, there exists one and only one $A$-homomorphism $w:C\to G$ such that $u=p_1\circ w$, $v=p_2\circ w$.
One can still say that $E\times_B F$ is the \emph{projective limit} of the projective system formed by $B$, $E$, $F$ and the $A$-homomorphisms $f$, $g$, in the category of $A$-rings ($\mathbf{0_{III}}$, 8.1.9).
\end{env}

\begin{env}[18.1.3]
\label{0.18.1.3}
Let $\mathfrak{J}$ be the bilateral ideal of $E$, kernel of $f$; it is immediate that the kernel $\mathfrak{J}'$ of $p_2:G\to F$ is the ideal formed of the elements $(x,0)$ where $x\in\mathfrak{J}$; the restriction $i_1:\mathfrak{J}'\to\mathfrak{J}$ of $p_1$ is an isomorphism of $A$-bimodules, and also an isomorphism of \unsure{$A$}-bimodules.
Similarly, if $\mathfrak{R}$ is the kernel of $g$, the kernel $\mathfrak{R}'$ of $p_1:G\to E$ is the ideal of the elements $(0,y)$ where $y\in\mathfrak{R}$, and the restriction $i_2:\mathfrak{R}'\to\mathfrak{R}$ of $p_2$ is an isomorphism (in the same sense).
Finally, it is clear that the kernel of the $A$-homomorphism $q=f\circ p_1=g\circ p_2$ of $E\times_B F$ into $B$ is $\mathfrak{J}\times\mathfrak{R}=\mathfrak{J}'\oplus\mathfrak{R}'$, so that one has a commutative diagram
\begin{equation}
\tag{18.1.3.1}
  \label{0.18.1.3.1}
  \xymatrix{
    & 0 & 0 & 0 &\\
    & \mathfrak{J}'\oplus\mathfrak{R}'\ar[u]\ar@{->>}[d] & & & \\
    0\ar[r] & \mathfrak{J}'\ar[r]\ar@{^{(}->}[u] & G\ar[r]^{p_2} & F\ar[r] & 0\\
    & & 0\ar[r] & \mathfrak{J}\ar@{^{(}->}[u] & \\
    0\ar[r] & \mathfrak{J}\ar[r] & E\ar[r]^f & B\ar[r] & 0
  }
\end{equation}

The definitions and results of \sref{0.18.1.2} and \sref{0.18.1.3} extend immediately to the fibre product of any \emph{family} $(E_\lambda)_{\lambda\in I}$ of $A$-rings defined by a family of $A$-homomorphisms $f_\lambda:E_\lambda\to B$.
We leave the formulation of these results to the reader.
\end{env}

\begin{env}[18.1.4]
\label{0.18.1.4}
Following the terminology of (M, VIII), we call an $A$-ring $E$ equipped with a \emph{surjective} $A$-homomorphism (called the \emph{augmentation} of $E$), $f:E\to B$, an \emph{$A$-ring augmented over $B$}; the kernel $\mathfrak{J}$ of $f$ is called the \emph{augmentation ideal}.
One says that the $A$-augmented ring $E$ is \emph{trivial} if there exists an $A$-homomorphism $s:B\to E$ that is \emph{right inverse} of the augmentation $f:E\to B$ (in other words $f\circ s=1_B$).
The exact sequence of $A$-bimodules
\[
  0\to\mathfrak{J}\to E\xrightarrow{f} B\to 0
\]
\oldpage[0\textsubscript{IV-1}]{53}
is then \emph{split}; in other words, one can identify the $A$-bimodule $E$ with $B\times\mathfrak{J}$, and the multiplication in $E$ is then given by
\[
  (b,z)(b',z')=(bb',bz'+zb'+zz'),
\]
$\mathfrak{J}$ being considered as a $B$-bimodule by means of $s:B\to E$.
\end{env}

\begin{env}[18.1.5]
\label{0.18.1.5}
With the notations of \sref{0.18.1.2}, suppose that the $A$-homomorphism $f$ makes $E$ into an $A$-ring \emph{augmented over $B$}.
Then it is clear that $p_2:G\to F$ is also surjective, in other words $G$ is given a structure of $A$-ring \emph{augmented over $F$}, that one calls the \emph{inverse image} by $g:F\to B$ of the augmented $A$-ring $E$.
\end{env}

\begin{proposition}[18.1.6]
\label{0.18.1.6}
For the inverse image $G=E\times_B F$ by $g:F\to B$ of the augmented $A$-ring $E$ to be a trivial augmented $A$-ring, it is necessary and sufficient that there exists an $A$-homomorphism $u:F\to E$ making the diagram
\[
  \xymatrix{
    & F\ar[ld]_u\ar[d]^g\\
    E\ar[r]_f & B
  }
\]
commutative.
\end{proposition}

\begin{proof}
The condition is evidently necessary, taking $u=p_1\circ s$, where $s:F\to G$ is an $A$-homomorphism right inverse of $p_2$.
Conversely, if there exists such an $A$-homomorphism $u$, the existence of the right inverse $s$ of $p_2$ follows from the universal property of the fibre product \sref{0.18.1.2} applied to the $A$-homomorphisms $u:F\to E$ and $1_F:F\to F$.

This result implies in particular that if $E$ is a trivial augmented $A$-ring, so is every inverse image.
\end{proof}

\begin{env}[18.1.7]
\label{0.18.1.7}
Let us take up again the situation described in \sref{0.18.1.2} and \sref{0.18.1.3}; on $\mathfrak{R}$ one has evidently a structure of $G$-\emph{bimodule}, coming from its structure of $F$-bimodule and the homomorphism of rings $p_2:G\to F$.
As $i_2$ is bijective one has moreover an injection $\theta=j\circ i_2^{-1}:\mathfrak{R}\to G$, which is a homomorphism of $G$-\emph{bimodules}.
We shall now see conversely that, when one supposes moreover that $f$ and $g$ are \emph{surjective}, in other words that $E$ and $F$ are $A$-rings \emph{augmented} over $B$, the data of such a homomorphism $\theta$ allows us to reconstruct the augmented ring $E$ (over $B$) from the augmented ring $G$ (over $F$).

More precisely, suppose we are given an $A$-ring $F$ augmented over $B$ and an $A$-ring $G$ augmented over $F$, the augmentation ideals being denoted $\mathfrak{R}$ and $\mathfrak{J}'$ respectively:
\begin{equation}
\tag{18.1.7.1}
  \label{0.18.1.7.1}
  \xymatrix{
    & 0\ar[d] & & \\
    & \mathfrak{R}\ar@/^/[ld]^\theta\ar[d] & & \\
    0\ar[r] & \mathfrak{J}'\ar[r] & G\ar[r]^h & F\ar[r] & 0\\
    & & & B\ar[u]_{g'}\ar[d] & \\
    & & & 0 &
  }
\end{equation}

The homomorphism $h$ defines on all the ideals of $F$ (and in particular on $\mathfrak{R}$) a structure of $G$-bimodule.
Let $\theta:\mathfrak{R}\to G$ be a homomorphism of $G$-bimodules making the diagram \sref{0.18.1.7.1} commutative; this implies that $\theta$ is \emph{injective} and that $\mathfrak{J}'\cap\theta(\mathfrak{R})=0$; as $\theta(\mathfrak{R})=\mathfrak{R}$, the composite homomorphism $g\circ h:G\to B$ is surjective; moreover the image $\mathfrak{J}$ of $\mathfrak{J}'$ in $E=G/\theta(\mathfrak{R})$, and the restriction of $q$ to $\mathfrak{J}'$ the canonical homomorphism $q:G\to E$ is injective.
One can therefore form the \emph{fibre product} $G'=E\times_B F$, and as the two $A$-homomorphisms $q:G\to E$ and $h:G\to F$ are such that $f\circ q=g\circ h$, they define a unique $A$-homomorphism $u:G'\to G$ by the universal property of $G'$ \sref{0.18.1.2}; we shall now see that $u$ is \emph{bijective}.
It suffices to show that $u$ is compatible with the finite filtrations on $G'$ and $G$ formed respectively by $G'$ and $\mathfrak{J}'=\operatorname{Ker}(G'\to F)$, and by $G$ and $\mathfrak{J}'$; moreover, one has seen \sref{0.18.1.3} that $\operatorname{gr}_0 u:F\to F$ is the identity and the assertion (ii) in turn follows from $\operatorname{gr}_1 u:\mathfrak{J}''\to\mathfrak{J}'$ is bijective, hence $u$ itself is bijective (Bourbaki, \emph{Alg.\ comm.}, chap.~III, \S~2, no.~8, cor.~3 of th.~1).
\end{env}

\subsection{Extensions of a ring by a bimodule}
\label{subsection:0.18.2}

\begin{env}[18.2.1]
\label{0.18.2.1}
Let $E$ be an $A$-ring augmented over $B$, $f:E\to B$ the augmentation, $\mathfrak{J}=\operatorname{Ker}(f)$ the augmentation ideal.
If $\mathfrak{J}^2=0$, $\mathfrak{J}$ is not only an $E$-bimodule but also a $B$-\emph{bimodule} since $B$ is isomorphic to $E/\mathfrak{J}$; more precisely, every $b\in B$ is of the form $f(x)$ with $x\in E$, and if $z$, $z'$ are in $\mathfrak{J}$, one has $(x+z)z'=xz'$ (resp.\ $z'(x+z)=z'x$) so that the value of $xz'$ (resp.\ $z'x$) does not depend on the element $x\in f^{-1}(b)$, and one can write $bz'$ (resp.\ $z'b$) which defines the structure of $B$-bimodule considered.
Conversely, if $\mathfrak{J}$ is equipped with a structure of $B$-bimodule such that $xz'=f(x)z'$ and $z'x=z'f(x)$ for $x\in E$ and $z'\in\mathfrak{J}$, it is clear that $\mathfrak{J}^2=0$ for this bimodule structure.
\end{env}

\begin{env}[18.2.2]
\label{0.18.2.2}
One calls \emph{$A$-extension of an $A$-ring $B$ by a $B$-bimodule $L$} an \emph{exact} sequence of homomorphisms of $A$-bimodules
\[
  0\to L\xrightarrow{j} E\xrightarrow{f} B\to 0
\]
where $E$ is an $A$-ring, $f$ an $A$-homomorphism of rings and one has, for $x\in E$ and $z\in L$,
\[
  j(f(x)z)=xj(z),\quad j(zf(x))=j(z)x
\]
hence \sref{0.18.2.1} $j(L)$ is a bilateral ideal of square zero of $E$.
By abuse of language one also says that $E$ is an extension of $B$ by $L$.
One says that two $A$-extensions $E$, $E'$ of $B$ by $L$ are \emph{$A$-equivalent} if there exists an isomorphism of $A$-rings $u:E\xrightarrow{\sim} E'$ making the diagram
\begin{equation}
\tag{18.2.2.1}
  \label{0.18.2.2.1}
  \xymatrix{
    0\ar[r] & L\ar@/^/[r]\ar@/_/[r] & E\ar[d]^u\ar@/^/[r] & B\ar[r] & 0\\
    0\ar[r] & L\ar@/^/[r]\ar@/_/[r] & E'\ar@/^/[r] & B\ar[r] & 0
  }
\end{equation}
commutative.
\end{env}

\begin{env}[18.2.3]
\label{0.18.2.3}
\oldpage[0\textsubscript{IV-1}]{55}
One says that an $A$-extension $E$ of $B$ by a $B$-bimodule $L$ is \emph{$A$-trivial} if $E$ is a trivial augmented $A$-ring (over $B$) \sref{0.18.1.4}.
One defines on the $A$-bimodule product $B\times L$ a structure of $A$-ring by setting $(x,s)(y,t)=(xy,xt+sy)$, as one easily checks, and the canonical maps $j:L\to B\times L$ and $f:B\times L\to B$ define an $A$-extension of $B$ by $L$ that is $A$-trivial, the canonical map $g:B\to B\times L$ being an $A$-homomorphism right inverse of $f$.
One says this extension is the \emph{trivial type extension} of $B$ by $L$ and one denotes it $D_B(L)$; it is immediate that every $A$-trivial $A$-extension of $B$ by a bimodule is $A$-equivalent to $D_B(L)$.

One will note that every extension of the $A$-ring $A$ itself by a bimodule is necessarily $A$-trivial.
\end{env}

\begin{env}[18.2.4]
\label{0.18.2.4}
Given two $A$-extensions $L\xrightarrow{j} E\xrightarrow{f} B$, $L'\xrightarrow{j'} E'\xrightarrow{f'} B'$, a \emph{morphism} of the first into the second is by definition a triplet of homomorphisms of $A$-bimodules $(u,v,w)$ such that the diagram
\begin{equation}
\tag{18.2.4.1}
  \label{0.18.2.4.1}
  \xymatrix{
    0\ar[r] & L\ar[r]^j\ar[d]_w & E\ar[r]^f\ar[d]_u & B\ar[r]\ar[d]_v & 0\\
    0\ar[r] & L'\ar[r]_{j'} & E'\ar[r]_{f'} & B'\ar[r] & 0
  }
\end{equation}
is commutative, $u$ and $v$ being $A$-homomorphisms of \emph{rings} and $w$ being such that $w(bz)=v(b)w(z)$ and $w(zb)=w(z)v(b)$ for $z\in L$ and $b\in B$ (in other words, the pair $(v,w)$ constitutes a \emph{di-homomorphism} of the $B$-bimodule $L$ into the $B'$-bimodule $L'$);
it is clear that if $(u',v',w')$ is a morphism from $L'\to E'\to B'$ into an $A$-extension $L''\to E''\to B''$, $(u'\circ u,v'\circ v,w'\circ w)$ is still a morphism, which justifies the terminology.

The consideration of the two commutative squares of the diagram \sref{0.18.2.4.1} will lead us to two operations on extensions of $A$-rings.
\end{env}

\begin{env}[18.2.5]
\label{0.18.2.5}
In the first place, consider an $A$-extension $E'$ of $B'$ by $L'$
\[
  0\to L'\xrightarrow{j'} E'\xrightarrow{f'} B'\to 0
\]
and an $A$-homomorphism of rings $v:B\to B'$, and let $F=E'\times_{B'} B$ be the \emph{inverse image} of $E'$ \sref{0.18.1.5}, so that one has a commutative diagram
\[
  \xymatrix{
    0\ar[r] & L_0\ar[r]\ar[d]^{i} & F\ar[r]^{p_2}\ar[d]^{p_1} & B\ar[r]\ar[d]^v & 0\\
    0\ar[r] & L'\ar[r]_{j'} & E'\ar[r]_{f'} & B'\ar[r] & 0
  }
\]
where the two lines are exact, $p_1$ and $p_2$ are canonical homomorphisms and $i$ is bijective; one sees from the definition \sref{0.18.1.2} and \sref{0.18.1.3} that $L_0^2=0$, so that one can consider $F$ as an $A$-extension of $B$ by $L'$ (where $L'$ being naturally considered as a $B$-bimodule by means of the homomorphism of rings $v$).
The functorial character of the fibre product with respect to each of the factors shows moreover that if one has a morphism of two extensions of $B'$
\[
  \xymatrix{
    0\ar[r] & L\ar[r]\ar[d]^{1_L} & E_1\ar[r]\ar[d]^g & B_1\ar[r]\ar[d]^h & 0\\
    0\ar[r] & L\ar[r] & E_2\ar[r] & B_2\ar[r] & 0
  }
\]
one deduces from it canonically a morphism of inverse image extensions
\[
  \xymatrix{
    0\ar[r] & L'\ar[r]\ar[d]^{1_{L'}} & E_1\oplus_L L'\ar[r]\ar[d]^{g\oplus 1_{L'}} & B_1\ar[r]\ar[d]^h & 0\\
    0\ar[r] & L'\ar[r] & E_2\oplus_L L'\ar[r] & B_2\ar[r] & 0
  }
\]

In particular, if $E_1$ and $E_2$ are $A$-equivalent $A$-extensions of $B$ by $L$, their inverse image extensions of $B$ by $L'$ obtained by $w$ are $A$-equivalent.

When a morphism \sref{0.18.2.4.1} of $A$-extensions factors through the inverse image of $E'$ by the homomorphism $w:L\to L'$ (where $L'$ is considered as $B$-bimodule by means of the homomorphism of rings $v:B\to B'$):
in fact, the definition of the amalgamated sum shows that there exists a unique $A$-homomorphism $u_0:E\to F=E'\times_{B'} B$ making the diagram
\[
  \xymatrix{
    0\ar[r] & L\ar[r]^j\ar[d]^{w_0} & E\ar[r]^f\ar[d]^{u_0} & B\ar[r]\ar[d]^{1_B} & 0\\
    0\ar[r] & L'_0\ar[r]\ar[d]^{i} & F\ar[r]^{p_2}\ar[d]^{p_1} & B\ar[r]\ar[d]^v & 0\\
    0\ar[r] & L'\ar[r]_{j'} & E'\ar[r]_{f'} & B'\ar[r] & 0
  }
\]
commutative, with $u=u_0\circ j_1$, $w=j_0\circ w_0$, $j_1$ and $j_2$ being the canonical homomorphisms; one verifies immediately that $u_0$ is also a homomorphism of rings.
\end{env}

\begin{env}[18.2.6]
\label{0.18.2.6}
Let us study the inverse image of extensions by surjective $A$-homomorphisms.
Consider a surjective $A$-homomorphism $v:B\to B'$, an $A$-extension $F$ of $B$ by a $B$-bimodule $L$, and the bilateral ideal $\mathfrak{R}$ of $B$, kernel of $v$ (which can be considered as an $F$-bimodule by means of the augmentation $F\to B$); one has seen \sref{0.18.1.7} that every homomorphism of $F$-bimodules $\theta:\mathfrak{R}\to F$ making the diagram
\[
  \xymatrix{
    & \mathfrak{R}\ar@/^/[ld]^\theta\ar[d]\\
    F\ar[r] & B
  }
\]
commutative determines an extension $E'$ of $B'=B/\mathfrak{R}$ by $L$ whose inverse image by $v:B\to B'$ is $A$-equivalent to $F$, and that every $A$-extension $E$ of $B'$ by $L$ having this last property is obtained in this way (up to $A$-equivalence).
Moreover, for two homomorphisms $\theta_1$, $\theta_2$ of $F$-bimodules of $\mathfrak{R}$ into $F$ to give two $A$-equivalent $A$-extensions of $B$ by $L$, it is necessary and sufficient that there exists an $A$-\emph{equivalence} $u$ of the $A$-extension $F$ onto itself such that $\theta_2=u\circ\theta_1$.
\end{env}

\begin{env}[18.2.7]
\label{0.18.2.7}
Consider now the left square of \sref{0.18.2.4.1}, and recall first the notion of \emph{amalgamated sum} in the category of $A$-bimodules: given three $A$-bimodules $X$, $Y$, $Z$ and two $A$-homomorphisms $f:X\to Y$, $g:X\to Z$, the amalgamated sum $Y\oplus_X Z$ is the \emph{inductive limit} of the inductive system formed by $X$, $Y$, $Z$ and the $A$-homomorphisms $f$, $g$ in the category of $A$-bimodules ($\mathbf{0_{III}}$, 8.1.11).
One defines this $A$-bimodule as quotient of the product $Y\times Z$ by the sub-$A$-bimodule $M$ image of $X$ by the homomorphism $x\mapsto(f(x),-g(x))$.
Its characteristic property is that, for every pair of homomorphisms of $A$-bimodules $u:Y\to T$, $v:Z\to T$ such that $u\circ f=v\circ g$, there exists a unique homomorphism $w:Y\oplus_X Z\to T$ such that $u=w\circ j_1$ and $v=w\circ j_2$, where $j_1:Y\to Y\oplus_X Z$ and $j_2:Z\to Y\oplus_X Z$ are the canonical maps.
\end{env}

\begin{env}[18.2.8]
\label{0.18.2.8}
Consider now an $A$-extension $E$ of $B$ by a $B$-bimodule $L$:
\[
  0\to L\xrightarrow{j} E\xrightarrow{f} B\to 0
\]
and let moreover $w:L\to L'$ be a homomorphism of $B$-\emph{bimodules}.
Let $H$ be the $A$-bimodule amalgamated sum $E\oplus_L L'$; let us show how one can equip it with a structure of $A$-\emph{ring} and define an $A$-\emph{extension}
\[
  0\to L'\to H\to B\to 0.
\]

Let us note for this that $L'$ is equipped with a structure of $E$-\emph{bimodule} by means of the augmentation homomorphism $E\to B$; one can therefore form the $A$-extension of trivial type $G=D_B(L')$ \sref{0.18.2.3}.
Consider then the map $\theta:z\mapsto(j(z),-w(z))$ from $L$ into $G$; it is a homomorphism of $G$-\emph{bimodules} ($L$ being considered as a $G$-bimodule by means of the canonical homomorphism $\rho:G\to E$).
Indeed, for $(x,z')\in G$ and $z\in L$, one has $(j(z),-w(z))(x,z')=(j(z)x,j(z)z'-w(z)f(x))$; but $j(z)z'=f(j(z))z'=0$ by definition of the structure of $E$-bimodule on $L'$, and $w(z)f(x)=w(zf(x))$ by definition of the structure of $B$-bimodule on $L'$; one thus verifies that $\theta$ is also a homomorphism of $G$-modules on the left.
One can then apply the result of \sref{0.18.1.7} to the diagram
\[
  \xymatrix{
    & 0\ar[d]\\
    & L\ar[d]\ar@/^/[ld]^\theta\\
    0\ar[r] & L'\ar[r] & G\ar[r]^p & E\ar[r] & 0\\
    & & & B\ar[u]_f\ar[d]\\
    & & & 0
  }
\]
the result of \sref{0.18.1.7}.
As $H=G/\theta(L)$ by definition, our assertion is an immediate consequence of \sref{0.18.1.7}.

One says that the $A$-extension $H$ of $B$ by $L'$ is \emph{deduced} of $E$ \emph{by the homomorphism} $w:L\to L'$.
The functorial character of the amalgamated sum in each of the summands shows moreover that if one has a morphism of extensions
\[
  \xymatrix{
    0\ar[r] & L\ar[r]\ar[d]^{1_L} & E_1\ar[r]\ar[d]^g & B_1\ar[r]\ar[d]^h & 0\\
    0\ar[r] & L\ar[r] & E_2\ar[r] & B_2\ar[r] & 0
  }
\]
one deduces from it canonically a morphism of deduced extensions
\[
  \xymatrix{
    0\ar[r] & L'\ar[r]\ar[d]^{1_{L'}} & E_1\oplus_L L'\ar[r]\ar[d] & B_1\ar[r]\ar[d]^h & 0\\
    0\ar[r] & L'\ar[r] & E_2\oplus_L L'\ar[r] & B_2\ar[r] & 0
  }
\]

In particular, if $E_1$ and $E_2$ are $A$-equivalent $A$-extensions of $B$ by $L$, the extensions of $B$ by $L'$ obtained by $w$ are $A$-equivalent.

When one has a morphism \sref{0.18.2.4.1} of $A$-extensions, it factors through the $A$-extension $H$ of $B$ by $L'$ deduced from $E$ by the homomorphism $w:L\to L'$ (where $L'$ being considered as $B$-bimodule by means of the homomorphism of rings $v:B\to B'$):
in fact, the definition of the amalgamated sum shows that there exists a unique $A$-homomorphism $u_0:H\to E'$ of $A$-bimodules, making the diagram
\[
  \xymatrix{
    0\ar[r] & L\ar[r]^j\ar[d]^{w_0} & E\ar[r]^f\ar[d]^{j_1} & B\ar[r]\ar[d]^{1_B} & 0\\
    0\ar[r] & L'_0\ar[r]\ar[d]^{j_e} & H\ar[r]^{f_e}\ar[d]^{u_e} & B\ar[r]\ar[d]^v & 0\\
    0\ar[r] & L'\ar[r]_{j'} & E'\ar[r]_{f'} & B'\ar[r] & 0
  }
\]
commutative, with $u=u_0\circ j_1$, $w=j_0\circ w_0$, $j_1$ and $j_2$ being the canonical homomorphisms; one verifies immediately that $u_0$ is also a homomorphism of rings.

Let us note moreover the functorial properties relative to trivial extensions:
\end{env}

\begin{proposition}[18.2.9]
\label{0.18.2.9}
Let $B$, $B'$ be two $A$-rings, $L$ a $B$-bimodule, $L'$ a $B'$-bimodule, $v:B\to B'$ an $A$-homomorphism of rings, $w:L\to L'$ a homomorphism of $A$-bimodules such that $(v,w)$ is a di-homomorphism of bimodules.
Then there exists a unique $A$-homomorphism of rings $u:D_B(L)\to D_{B'}(L')$ making commutative the diagrams
\[
  \xymatrix{
    D_B(L)\ar[r]^u & D_{B'}(L') & & D_B(L)\ar[r]^u & D_{B'}(L')\\
    B\ar[u]\ar[r]_v & B'\ar[u] & & L\ar[u]\ar[r]_w & L'\ar[u]
  }
\]
where the vertical arrows are the canonical injections.
\end{proposition}

\begin{proof}
Indeed, $u$ cannot be other than the map $(x,s)\mapsto(v(x),w(s))$ and it remains to verify that it is an $A$-homomorphism of rings, which follows trivially from the definition of the multiplication \sref{0.18.2.3}.
\end{proof}

\begin{proposition}[18.2.10]
\label{0.18.2.10}
Let $B$ be an $A$-ring, $L$ a $B$-bimodule, $E$ an $A$-extension of $B$ by $L$.
One defines a bijective map of the set $G$ of homomorphisms of $B$-rings of $E$ onto the set $G'$ of $A$-equivalences of $D_B(L)$ onto $E$ by associating to every $g\in G$ the $A$-equivalence $g':(x,s)\mapsto g(x)+s$; the inverse map associates to every $g'\in G'$ the $B$-homomorphism right inverse of the augmentation $E\to B$.
\end{proposition}

\begin{proof}
This follows immediately from the definitions.
\end{proof}

\subsection{The group of classes of A-extensions}
\label{subsection:0.18.3}

\begin{env}[18.3.1]
\label{0.18.3.1}
Consider an $A$-ring $B$ and a $B$-bimodule $L$ fixed; then the relation ``$E$ and $E'$ are $A$-equivalent'' between $A$-extensions $E$, $E'$ of $B$ by $L$ is an equivalence relation and for this relation one can speak of the \emph{set} of classes of $A$-extensions of $B$ by $L$.
To see this, it suffices to remark that, for every $x\in B$, if $c_x$ is an element of $E$ whose image in $B$ is $x$, every $z\in E$ can be written in one and only one way in the form $c_x+t$, where $t\in L$, and one can write $c_{x+y}=c_x+c_y+\varphi(x,y)$, $c_{xy}=c_x c_y+\psi(x,y)$, where $\varphi(x,y)$ and $\psi(x,y)$ are elements of $L$, the maps $\varphi$ and $\psi$ from $B\times B$ into $L$ being the expressions of the fact that $E$ is an $A$-ring, which we shall not write out here.
Every $A$-extension of $B$ by $L$ is therefore $A$-equivalent to an $A$-extension whose underlying set is $B\times L$, hence our conclusion.
\end{env}

\begin{env}[18.3.2]
\label{0.18.3.2}
The $A$-ring $B$ being fixed, denote provisionally by $T(L)$ the set of classes of $A$-extensions of $B$ by $L$.
For every $B$-homomorphism of $B$-bimodules $w:L\to L'$, one defines canonically a map $T(w):T(L)\to T(L')$ by associating to the class of an $A$-extension $E$ of $B$ by $L$ the class of the $A$-extension $E\oplus_L L'$ deduced from $E$ by $w$, by virtue of \sref{0.18.2.8}.
If $w':L'\to L''$ is a second homomorphism of $B$-bimodules, one has also
\begin{equation}
\tag{18.3.2.1}
  \label{0.18.3.2.1}
  T(w'\circ w)=T(w')\circ T(w).
\end{equation}

Indeed, one knows that there exists a canonical isomorphism of $A$-bimodules
\[
  E\oplus_L L''\xrightarrow{\sim}(E\oplus_L L')\oplus_{L'} L''
\]
by virtue of the general properties of inductive limits (cf.\ for example ($\mathbf{I}$, 3.3.9)), and it is immediate to verify that it is an $A$-equivalence of $A$-extensions, hence our assertion.
One will identify canonically $T(L)$ to $\prod_\alpha T(L_\alpha)$ by the preceding map.
One verifies also immediately that if $(L'_\alpha)_{\alpha\in I}$ is a second family of $B$-bimodules and, for each $\alpha$, $w_\alpha:L_\alpha\to L'_\alpha$ a homomorphism of $B$-bimodules, then, setting $w=\prod_\alpha w_\alpha:L\to L'=\prod_\alpha L'_\alpha$, $T(w)$ identifies with $\prod_\alpha T(w_\alpha)$ when one makes the preceding identification.
\end{env}

\begin{env}[18.3.3]
\label{0.18.3.3}
Consider now a family $(L_\alpha)_{\alpha\in I}$ of $B$-bimodules and their product $L=\prod_\alpha L_\alpha$; the projections $\operatorname{pr}_\alpha:L\to L_\alpha$ define applications
\[
  T(\operatorname{pr}_\alpha):T(L)\to T(L_\alpha),
\]
hence a canonical map
\begin{equation}
\tag{18.3.3.1}
  \label{0.18.3.3.1}
  \prod_\alpha T(\operatorname{pr}_\alpha):T(L)\to\prod_{\alpha\in I} T(L_\alpha).
\end{equation}

We shall see that this map is \emph{bijective}.
Indeed, for every $\alpha\in I$, let $E_\alpha$ be an $A$-extension of $B$ by $L_\alpha$, and let $F$ be the \emph{fibre product} of the $E_\alpha$ on $B$ \sref{0.18.1.3}; it is immediate that if one replaces each $E_\alpha$ by an $A$-equivalent $A$-extension $E'_\alpha$, the fibre product $F'$ of the $E'_\alpha$ on $B$ is $A$-equivalent to $F$.
One has thus defined a map $\prod_\alpha T(L_\alpha)\to T(L)$ and it is clear that this map is the inverse of \sref{0.18.3.3.1}, hence our assertion.
One will identify canonically $T(L)$ with $\prod_\alpha T(L_\alpha)$ by the preceding map.
One verifies also immediately that if $(L'_\alpha)_{\alpha\in I}$ is a second family of $B$-bimodules and, for each $\alpha$, $w_\alpha:L_\alpha\to L'_\alpha$ a homomorphism of $B$-bimodules, then, setting $w=\prod_\alpha w_\alpha:L\to L'=\prod_\alpha L'_\alpha$, $T(w)$ identifies with $\prod_\alpha T(w_\alpha)$ when one makes the preceding identification.
\end{env}

\begin{env}[18.3.4]
\label{0.18.3.4}
For a $B$-bimodule $L$, the addition is a homomorphism $s:L\times L\to L$ of $B$-\emph{bimodules}, and so is the symmetry $t:L\to L$ of the additive law of $L$.
One deduces from them a law of composition
\[
  T(s):T(L)\times T(L)\to T(L)
\]
in $T(L)$ by virtue of \sref{0.18.3.3}, and this law is a law of \emph{commutative group} of which $T(t)$ is the symmetry, as it follows from the definition of an object in groups by means of commutative diagrams ($\mathbf{0_{III}}$, 8.2.5 and 8.2.6).
We shall denote by $\operatorname{Exan}_A(B,L)$ the commutative group thus defined and we say it is the \emph{group of classes of $A$-extensions of $B$ by $L$}.
\end{env}

\begin{env}[18.3.5]
\label{0.18.3.5}
\oldpage[0\textsubscript{IV-1}]{61}
Denote by $\boldsymbol{K}$ the category whose objects are the triplets $(A,B,L)$ where $A$ is a ring, $B$ an $A$-ring and $L$ a $B$-bimodule; the \emph{morphisms} of this category are the triplets $(u,v,w)$ where $u:A'\to A$ and $v:B'\to B$ are two homomorphisms of rings making commutative the left square of the diagram
\[
  \xymatrix{
    A\ar[r] & B & L\\
    A'\ar[u]^u\ar[r] & B'\ar[u]^v & L'\ar[u]_w
  }
\]
where the horizontal arrows are the structural homomorphisms; finally $w:L\to L'$ is a homomorphism of commutative groups such that one has $w(v(b')z)=b'w(z)$ and $w(zv(b'))=w(z)b'$ for all $z\in L$ and $b'\in B'$ (in other words, $w$ is a homomorphism of $B'$-bimodules when $L$ is equipped with its structure of $B'$-bimodule defined by $v$).
The composition of morphisms is defined by $(u',v',w')\circ(u,v,w)=(u\circ u',v\circ v',w'\circ w)$, which is immediately verified.
We propose to show that
\begin{equation}
\tag{18.3.5.1}
  \label{0.18.3.5.1}
  (A,B,L)\mapsto\operatorname{Exan}_A(B,L)
\end{equation}
is a \emph{covariant functor} from the category $\boldsymbol{K}$ into the category $\boldsymbol{Ab}$ of commutative groups.
It is therefore a matter of defining for every triplet $(u,v,w)$ as above, a homomorphism of commutative groups
\[
  (u,v,w)_*:\operatorname{Exan}_A(B,L)\to\operatorname{Exan}_{A'}(B',L').
\]

By virtue of the definition of morphisms in $\boldsymbol{K}$, one can write
\[
  (u,v,w)=(1_{A'},1_{B'},w)\circ(1_{A'},v,1_L)\circ(u,1_B,1_L)
\]
where, in the first factor, $L$ is equipped with its structure of $B'$-bimodule defined by $v$; we shall therefore define $(u,v,w)_*$ when two of the three homomorphisms $u$, $v$, $w$ are reduced to the identity.
\end{env}

\begin{env}[18.3.6]
\label{0.18.3.6}
We take first $(1_A,1_B,w)_*$ the map
\begin{equation}
\tag{18.3.6.1}
  \label{0.18.3.6.1}
  w_*:\operatorname{Exan}_A(B,L)\to\operatorname{Exan}_A(B,L')
\end{equation}
denoted $T(w)$ in \sref{0.18.3.2}; it is immediate to verify that it is a homomorphism of groups, this property expressing the commutativity of diagrams, transformed by $T$ of analogous diagrams for $L$ and $L'$.

The map $(1_A,v,1_L)_*$ is the map
\begin{equation}
\tag{18.3.6.2}
  \label{0.18.3.6.2}
  v^*:\operatorname{Exan}_A(B,L)\to\operatorname{Exan}_A(B',L)
\end{equation}
defined in the following way: if $E$ is an $A$-extension of $B$ by $L$, one has seen that $E\times_B B'$ is an $A$-extension of $B'$ by $L$ \sref{0.18.2.5} and that if one replaces $E$ by an $A$-equivalent $A$-extension $E'$, $E'\times_B B'$ is $A$-equivalent to $E\times_B B'$; the image by $v^*$ of the class of $E$ is the class of $E\times_B B'$.
One verifies immediately that if $w:L\to L'$ is a homomorphism of $B$-bimodules, the diagram
\begin{equation}
\tag{18.3.6.3}
  \label{0.18.3.6.3}
  \xymatrix{
    \operatorname{Exan}_A(B,L)\ar[r]^{v^*}\ar[d]_{w_*} & \operatorname{Exan}_{A'}(B,L)\ar[d]^{w_*}\\
    \operatorname{Exan}_A(B,L')\ar[r]_{v^*} & \operatorname{Exan}_{A'}(B,L')
  }
\end{equation}
is commutative, $L$ and $L'$ being considered as $B'$-bimodules by means of $v$ in the right column.
Replacing $L$ and $L'$ by $L\times L$ and $L$, and $w$ by the addition $s$ in $L$, one concludes as above that $v^*$ is indeed a \emph{homomorphism of groups}.
\end{env}

\begin{env}[18.3.6.4]
\label{0.18.3.6.4}
Finally, the map $(u,1_B,1_L)_*$ is the map
\begin{equation}
  u^*:\operatorname{Exan}_A(B,L)\to\operatorname{Exan}_{A'}(B,L)
\end{equation}
which is obtained by associating to an $A$-extension $E$ of $B$ by $L$ the ring $E$ considered as an $A'$-\emph{ring} by means of $u$ \sref{0.18.1.1}, which is evidently also an $A'$-extension of $B$ by $L$, $B$ being also considered as $A'$-ring by means of $u$; it is clear that an $A$-equivalence is also an $A'$-equivalence, whence the map \sref{0.18.3.6.4}, which, for every homomorphism $w:L\to L'$ of $B$-bimodules, also makes commutative the diagram analogous to that
\[
  \xymatrix{
    \operatorname{Exan}_A(B,L)\ar[r]^{u^*}\ar[d]_{w_*} & \operatorname{Exan}_{A'}(B,L)\ar[d]^{w_*}\\
    \operatorname{Exan}_A(B,L')\ar[r]_{u^*} & \operatorname{Exan}_{A'}(B,L')
  }
\]
One deduces as above that $u^*$ is indeed a homomorphism of groups.

This being so, one sets $(u,v,w)_*=(1_{A'},1_{B'},w)_*\circ(1_{A'},v,1_L)_*\circ(u,1_B,1_L)_*$ and one verifies easily, by reason of the commutativity of the diagrams \sref{0.18.3.6.3} and \sref{0.18.3.6.5}, that one has $(u\circ u',v\circ v',w'\circ w)_*=(u',v',w')_*\circ(u,v,w)_*$, which finishes the proof that \sref{0.18.3.5.1} is a functor.

The existence of the homomorphism of groups \sref{0.18.3.6.1} shows in particular that if $E$ is a trivial $A$-extension of $B$ by $L$, the extension $E\oplus_L L'$ of $B$ by $L'$ defined in \sref{0.18.2.8} is also trivial, which one moreover verifies directly without difficulty.

On the other hand, for every element $z$ of the \emph{centre} $Z$ of $B$, the homothety $h_z:y\mapsto yz=zy$ is an endomorphism of the $B$-bimodule $L$, hence $(h_z)_*$ is an endomorphism of the commutative group $\operatorname{Exan}_A(B,L)$, and by functoriality, these endomorphisms define on $\operatorname{Exan}_A(B,L)$ a structure of $Z$-\emph{module}.
\end{env}

\begin{env}[18.3.7]
\label{0.18.3.7}
Let $A$, $A'$ be two rings, $u:A'\to A$ a homomorphism, $B$ an $A$-ring and $L$ a $B$-bimodule.
The \emph{kernel} of the homomorphism of groups
\[
  u^*:\operatorname{Exan}_A(B,L)\to\operatorname{Exan}_{A'}(B,L)
\]
is formed by definition of classes of $A$-extensions of $B$ by $L$ that are $A'$-trivial when one considers them as $A'$-extensions by means of $u$.
One denotes this kernel by the notation $\operatorname{Exan}_{A/A'}(B,L)$ when this causes no confusion.

If $\Lambda$ is a ring, and if one has a commutative diagram of $\Lambda$-\emph{homomorphisms} of $\Lambda$-rings
\begin{equation}
\tag{18.3.7.1}
  \label{0.18.3.7.1}
  \xymatrix{
    B'\ar[r] & B\\
    A'\ar[u]\ar[r] & A\ar[u]
  }
\end{equation}
\oldpage[0\textsubscript{IV-1}]{63}
one deduces from it canonically homomorphisms
\begin{equation}
\tag{18.3.7.2}
  \label{0.18.3.7.2}
  \operatorname{Exan}_{A/\Lambda}(B,L)\to\operatorname{Exan}_{A'/\Lambda}(B,L)\to\operatorname{Exan}_{A'/\Lambda}(B',L)
\end{equation}
which come from those of \sref{0.18.3.7.1} by functoriality.
\end{env}

\begin{proposition}[18.3.8]
\label{0.18.3.8}
Let $B$ be a ring, $\mathfrak{J}$ a bilateral ideal of $B$, $C=B/\mathfrak{J}$ the quotient ring canonically equipped with a structure of $C$-bimodule.
For every $C$-bimodule $L$, the group $\operatorname{Hom}_C(\mathfrak{J}/\mathfrak{J}^2,L)$ of the additive group of homomorphisms of $C$-bimodules of $\mathfrak{J}/\mathfrak{J}^2$ into $L$.
One defines then an isomorphism of commutative groups
\begin{equation}
\tag{18.3.8.1}
  \label{0.18.3.8.1}
  \eta_L:\operatorname{Hom}_C(\mathfrak{J}/\mathfrak{J}^2,L)\xrightarrow{\sim}\operatorname{Exan}_B(C,L)
\end{equation}
by associating to every $C$-homomorphism $w:\mathfrak{J}/\mathfrak{J}^2\to L$ (which is \emph{a fortiori} a $B$-homomorphism) the class of the extension $(B/\mathfrak{J}^2)\oplus_{(\mathfrak{J}/\mathfrak{J}^2)} L$ deduced from the extension $B/\mathfrak{J}^2$ of $C$ by $\mathfrak{J}/\mathfrak{J}^2$ by means of $w$; the inverse isomorphism associates to a class of $B$-extension $E$ of $C$ by $L$ the homomorphism $w$ such that the composite $\mathfrak{J}\to\mathfrak{J}/\mathfrak{J}^2\xrightarrow{w} L$ is the restriction to $\mathfrak{J}$ of the structural homomorphism $B\to E$.

Let $F=B/\mathfrak{J}^2$, which is a $B$-extension of $C$ by $\mathfrak{J}/\mathfrak{J}^2$.
For every $B$-extension
$0\to L\xrightarrow{j} E\xrightarrow{p} C\to 0$ of $C$ by $L$, the structural homomorphism $f:B\to E$ is such that the composite $B\xrightarrow{f} E\xrightarrow{p} C$ is the canonical homomorphism $B\to B/\mathfrak{J}=C$.
As the image of $\mathfrak{J}$ by $p\circ f$ is zero, $f(\mathfrak{J})$ is contained in the kernel $j(L)$, and as $j(L)$ is of square zero, $f(\mathfrak{J}^2)=0$; hence $f$ factors as $B\to B/\mathfrak{J}^2=F\xrightarrow{u} E$, and if $w$ is the restriction of $u$ to $\mathfrak{J}/\mathfrak{J}^2$, one has a commutative diagram
\[
  \xymatrix{
    0\ar[r] & \mathfrak{J}/\mathfrak{J}^2\ar[r]\ar[d]^w & F\ar[r]^g\ar[d]^u & C\ar[r]\ar@{=}[d] & 0\\
    0\ar[r] & L\ar[r]_j & E\ar[r]_p & C\ar[r] & 0
  }
\]
in other words $(u,1_C,w)$ is a morphism of extensions \sref{0.18.2.4}.
One deduces from it a morphism of extensions
\[
  \xymatrix{
    0\ar[r] & L\ar[r]\ar@{=}[d] & E'\ar[r]\ar[d]^{u'} & C\ar[r]\ar@{=}[d] & 0\\
    0\ar[r] & L\ar[r]_j & E\ar[r]_p & C\ar[r] & 0
  }
\]
it remains to prove that $u'$ is a $B$-\emph{equivalence}, which establishes $\eta_L$ is bijective.
It suffices to prove \sref{0.18.1.7} that the inverse image $G$ of the extension $E$ of $C$ by $F$ is $F$-\emph{trivial}, which is evident since $g$ factors as $F\xrightarrow{u} E\xrightarrow{p} C$ \sref{0.18.1.6}.

It remains to see that $\eta_L$ is a homomorphism of \emph{groups}; now, for every $B$-homomorphism $h:L\to L'$, it is immediate that the diagram
\[
  \xymatrix{
    \operatorname{Hom}_C(\mathfrak{J}/\mathfrak{J}^2,L)\ar[r]^{\eta_L}\ar[d]^{\operatorname{Hom}(1,h)} & \operatorname{Exan}_B(C,L)\ar[d]^{h_*}\\
    \operatorname{Hom}_C(\mathfrak{J}/\mathfrak{J}^2,L')\ar[r]_{\eta_{L'}} & \operatorname{Exan}_B(C,L')
  }
\]
is commutative.
It suffices to apply this remark to the homomorphism $L\times L\to L$ defining the addition to conclude.
\end{proposition}

\subsection{Extensions of algebras}
\label{subsection:0.18.4}

\begin{env}[18.4.1]
\label{0.18.4.1}
\oldpage[0\textsubscript{IV-1}]{64}
Let $A$ be a \emph{commutative} ring.
The category of $A$-\emph{algebras} is then a full subcategory of that of $A$-rings, characterized by the fact that the structural homomorphisms $\rho:A\to B$ are \emph{central}.

If $B$ is an $A$-algebra, $L$ a $B$-bimodule, it is immediate that the $A$-extension of trivial type $D_B(L)$ is an $A$-algebra.
Similarly, in the construction of \sref{0.18.2.5} (resp.\ of \sref{0.18.2.8}), if $B$, $B'$ and $E$ (resp.\ $B$ and $E$) are $A$-algebras, it is the same for $E'\times_{B'} B$ (resp.\ for $E\oplus_L L'$).
Finally, if $B$ is an $A$-algebra and $E$ an $A$-extension of $B$ by a $B$-bimodule $L$ that is an $A$-algebra, every $A$-extension $E'$ of $B$ by $L$, $A$-\emph{equivalent} to $E$, is also an $A$-algebra.
One deduces then immediately from the definition \sref{0.18.3.4} that the classes of $A$-extensions equivalent to $A$-\emph{algebras} form a \emph{subgroup}, denoted $\operatorname{Exal}_A(B,L)$, of $\operatorname{Exan}_A(B,L)$.
Let $\boldsymbol{K}'$ be the full subcategory of $\boldsymbol{K}$ defined in \sref{0.18.3.5}, whose objects $(A,B,L)$ are such that $A$ is commutative and $B$ an $A$-algebra.
Then what precedes shows that
\[
  (A,B,L)\mapsto\operatorname{Exal}_A(B,L)
\]
is a \emph{covariant functor} from $\boldsymbol{K}'$ to $\boldsymbol{Ab}$.
The results of \sref{0.18.3.7} are unchanged when one replaces $\operatorname{Exan}$ everywhere by $\operatorname{Exal}$.
\end{env}

\begin{env}[18.4.2]
\label{0.18.4.2}
Suppose always $A$ commutative.
The remarks of \sref{0.18.4.1} remain valid when one replaces ``algebra'' by ``commutative algebra'' and ``bimodule'' by ``module''.
If $B$ is a commutative $A$-algebra and $L$ a $B$-module, the classes of $A$-extensions equivalent to $A$-rings that are commutative $A$-\emph{algebras} (which amounts to the same, commutative $A$-\emph{rings}) form a \emph{subgroup} of $\operatorname{Exal}_A(B,L)$, denoted $\operatorname{Exalcom}_A(B,L)$.
If $\boldsymbol{K}''$ is the full subcategory of $\boldsymbol{K}'$ where $A$ and $B$ are commutative and $L$ a $B$-module,
\[
  (A,B,L)\mapsto\operatorname{Exalcom}_A(B,L)
\]
is still a \emph{covariant functor} from $\boldsymbol{K}''$ to $\boldsymbol{Ab}$.
One can also replace $\operatorname{Exan}$ by $\operatorname{Exalcom}$ in \sref{0.18.3.7}.
Finally, if in \sref{0.18.3.8}, one supposes $B$ to be a commutative ring and $L$ a $C$-\emph{module}, the same reasoning gives a canonical isomorphism
\begin{equation}
\tag{18.4.2.1}
  \label{0.18.4.2.1}
  \operatorname{Hom}_C(\mathfrak{J}/\mathfrak{J}^2,L)\xrightarrow{\sim}\operatorname{Exalcom}_B(C,L)
\end{equation}
where the first member is the group of homomorphisms of $C$-\emph{modules}.
\end{env}

\begin{env}[18.4.3]
\label{0.18.4.3}
Let $A$ be a commutative ring.
An important case of extensions of $A$-algebras is formed of $A$-algebras $E$, extensions of an $A$-algebra $B$ by a $B$-bimodule $L$ such that $E$, as $A$-\emph{module}, is the \emph{trivial} extension of the $A$-module $B$ by the $A$-module $L$; in other words, the exact sequence of $A$-\emph{modules} $0\to L\to E\to B\to 0$ is \emph{split}.
One says then that $E$ is a \emph{Hochschild extension} of $B$ by $L$.
This will always be the case when $B$ is an $A$-module \emph{projective}, and in particular when $A$ is a commutative \emph{field}.
As an $A$-module, one can identify $E$ with $B\times L$, the multiplication in $E$ being given by $(x,s)(y,t)=(xy,xt+sy+f(x,y))$ with $f(x,y)\in L$.
If one writes that this multiplication defines on $B\times L$ a structure of $A$-algebra, one finds (M, XIV, 2) that $f$ must be an $A$-\emph{bilinear} map of $B\times B$ into $L$, such that
\begin{equation}
\tag{18.4.3.1}
  \label{0.18.4.3.1}
  f(xy,z)+f(x,y)z=xf(y,z)+f(x,yz)
\end{equation}
in other words, $f$ is a $2$-\emph{cocycle} on $B$ with values in $L$, in the sense of Hochschild; for the extension $E$ to be $A$-trivial, it is necessary and sufficient that
\begin{equation}
\tag{18.4.3.2}
  \label{0.18.4.3.2}
  f(x,y)=xg(y)-g(xy)+g(x)y
\end{equation}
where $g$ is an $A$-\emph{linear} map of $B$ into $L$, in other words $f$ must be a $2$-\emph{coboundary} in the sense of Hochschild.
One deduces from it that the classes of Hochschild extensions of $B$ by $L$ form a subgroup of $\operatorname{Exal}_A(B,L)$, isomorphic to the \emph{Hochschild cohomology group} $\mathrm{H}^2_A(B,L)$.

If $B$ is a commutative $A$-algebra, and $L$ a $B$-module, the commutative Hochschild extensions of $B$ by $L$ correspond to the \emph{symmetric} $2$-cocycles, i.e.\ those such that $f(y,x)=f(x,y)$.
The classes of commutative Hochschild extensions of $B$ by $L$ form therefore a subgroup of $\operatorname{Exalcom}_A(B,L)$, isomorphic to the subgroup of $\mathrm{H}^2_A(B,L)$ image of the group of symmetric $2$-cocycles, which we shall denote $\mathrm{H}^2_A(B,L)^s$.
\end{env}

\begin{env}[18.4.4]
\label{0.18.4.4}
Let us restrict to the case where $A$ and $B$ are \emph{commutative} and $L$ a $B$-module, and recall in this case the equivalent definition of the Hochschild cohomology groups $\mathrm{H}^i_A(B,L)$ (M, IX, 6).
One considers a complex $\mathrm{P}_\bullet=(\mathrm{P}_n)_{n\geqslant 0}$ of $B$-modules, where $\mathrm{P}_n=B^{\otimes(n+1)}=B\otimes_A B\otimes\cdots\otimes_A B$ ($n+1$ times), the structure of $B$-module being defined by $x(y_1\otimes y_2\otimes\cdots\otimes y_{n+1})=(xy_1)\otimes y_2\otimes\cdots\otimes y_{n+1}$; the derivation, of degree $-1$, $d_n:\mathrm{P}_n\to\mathrm{P}_{n-1}$, is defined by
\begin{align*}
  d_n(x_1\otimes x_2\otimes\cdots\otimes x_{n+1})&=(x_1 x_2)\otimes x_3\otimes\cdots\otimes x_{n+1}-x_1\otimes(x_2 x_3)\otimes\cdots\otimes x_{n+1}+\cdots\\
  &\quad +(-1)^{n-1} x_1\otimes x_2\otimes\cdots\otimes(x_n x_{n+1})+(-1)^n(x_{n+1} x_1)\otimes x_2\otimes\cdots\otimes x_n
\end{align*}
which is $B$-linear since $B$ is commutative.

A $2$-cocycle of this complex with values in $L$ is a $B$-linear map $h$ of $B\otimes_A B\otimes_A B$ into $L$ such that $h(d_3(x\otimes y\otimes z\otimes t))=0$; but as $h(x\otimes y\otimes z)=xh(1\otimes y\otimes z)$, the cocycle $h$ is determined by the $A$-bilinear map $(y,z)\mapsto f(y,z)=h(1\otimes y\otimes z)$ of $B$ into $L$; one obtains $h'(d_2(x\otimes y\otimes z))=0$, which gives back the condition \sref{0.18.4.3.1}.
Similarly, a $2$-coboundary will be a map of the form $x\otimes y\otimes z\mapsto h'(d_2(x\otimes y\otimes z))$, where $h':B\otimes_A B\to L$ is linear, and here $h'$ is determined by the $A$-linear map $g:y\mapsto h'(1\otimes y)$ of $B$ into $L$; one obtains then $h'(d_2(1\otimes y\otimes z))=xg(y)-g(xy)+yg(x)$, which gives back \sref{0.18.4.3.2}.
One proceeds similarly for the $1$-cocycles and one sees that one has
\begin{equation}
\tag{18.4.4.1}
  \label{0.18.4.4.1}
  \mathrm{H}^i_A(B,L)=\mathrm{H}^i(\operatorname{Hom}_B(\mathrm{P}_\bullet,L)).
\end{equation}
\end{env}

\begin{env}[18.4.5]
\label{0.18.4.5}
Under the conditions of \sref{0.18.4.4}, one can reinterpret in the same way the group $\mathrm{H}^2_A(B,L)^s$ \sref{0.18.4.3}.
For this, modify the complex $\mathrm{P}_\bullet$ in degree $3$ considering a new complex
\[
  \mathrm{P}'_\bullet:\mathrm{P}'_3\xrightarrow{d'_3}\mathrm{P}_2\xrightarrow{d_2}\mathrm{P}_1;
\]
we take $\mathrm{P}'_3=\mathrm{P}_3\oplus(B\otimes_A B\otimes_A B)$, $d'_3$ coinciding with $d_3$ on $\mathrm{P}_3$, and being given by
\[
  d'_3(x\otimes y\otimes z)=x\otimes y\otimes z-x\otimes z\otimes y.
\]

The relation $d_2 d'_3=0$ follows from the above with the commutativity of $B$.
With the notations introduced above, a $2$-cocycle of $\mathrm{P}'_\bullet$ corresponds now to an $A$-bilinear map $f$ of $B\times B$ into $L$ that is \emph{symmetric} and satisfies \sref{0.18.4.3.1}; consequently, one has $\mathrm{H}^2_A(B,L)^s=\mathrm{H}^2(\operatorname{Hom}_B(\mathrm{P}'_\bullet,L))$.
\end{env}

\begin{env}[18.4.6]
\label{0.18.4.6}
In the particular case where one considers a commutative \emph{field} $k$, an extension $K$ of $k$, one has $\operatorname{Ext}^1_K(M,L)=0$ whatever the $K$-vector spaces $L$, $M$ and consequently (M, VI, 3.3 $a$)), one has a canonical isomorphism
\begin{equation}
\tag{18.4.6.1}
  \label{0.18.4.6.1}
  \mathrm{H}^1_k(K,L)\simeq\operatorname{Hom}_K(\mathrm{H}_1(\mathrm{P}_\bullet),L)
\end{equation}
and similarly
\begin{equation}
\tag{18.4.6.2}
  \label{0.18.4.6.2}
  \mathrm{H}^2_k(K,L)^s\simeq\operatorname{Hom}_K(\mathrm{H}_2(\mathrm{P}'_\bullet),L).
\end{equation}
\end{env}

\subsection{Case of topological rings}
\label{subsection:0.18.5}

\begin{env}[18.5.1]
\label{0.18.5.1}
Let $B$ be a \emph{topological} ring whose topology is \emph{linear}, $\rho:A\to B$ a continuous homomorphism, $L$ a topological $B$-bimodule, and suppose that there exists an \emph{open} bilateral ideal $\mathfrak{R}_0$ of $B$ such that $\mathfrak{R}_0.L=L.\mathfrak{R}_0=0$, so that $L$ can be
\oldpage[0\textsubscript{IV-1}]{67}
considered as a $(B/\mathfrak{R})$-bimodule for every open bilateral ideal $\mathfrak{R}\subset\mathfrak{R}_0$.
Let then $\mathfrak{R}$ be an open bilateral ideal of $A$ and $\mathfrak{J}$ an open bilateral ideal of $B$ such that $\rho(\mathfrak{J})\subset\mathfrak{R}$, so that $B/\mathfrak{R}$ can be considered as an $(A/\mathfrak{J})$-ring discretely; the preceding remark proves that the group $\operatorname{Exan}_{A/\mathfrak{J}}(B/\mathfrak{R},L)$ is defined; moreover, if $\mathfrak{R}'\subset\mathfrak{R}$, $\mathfrak{J}'\subset\mathfrak{J}$ are two open bilateral ideals of $A$ and $B$ respectively such that $\rho(\mathfrak{J}')\subset\mathfrak{R}'$, one has by \sref{0.18.3.5.1} a canonical homomorphism
\begin{equation}
\tag{18.5.1.1}
  \label{0.18.5.1.1}
  \operatorname{Exan}_{A/\mathfrak{J}}(B/\mathfrak{R},L)\to\operatorname{Exan}_{A/\mathfrak{J}'}(B/\mathfrak{R}',L).
\end{equation}

The set of pairs of open bilateral ideals $(\mathfrak{J},\mathfrak{R})$ such that $\rho(\mathfrak{J})\subset\mathfrak{R}\subset\mathfrak{R}_0$ is evidently directed for the relation ``$\mathfrak{J}\supset\mathfrak{J}'$ and $\mathfrak{R}\supset\mathfrak{R}'$'' and the maps \sref{0.18.5.1.1} define an \emph{inductive system} of additive groups with this set as index set.
One sets, by abuse of notation (since it is no longer a matter of a group of extensions in bijective natural correspondence with a set of extensions)
\begin{equation}
\tag{18.5.1.2}
  \label{0.18.5.1.2}
  \operatorname{Exantop}_A(B,L)=\varinjlim\operatorname{Exan}_{A/\mathfrak{J}}(B/\mathfrak{R},L).
\end{equation}

To say that the second member of \sref{0.18.5.1.2} is \emph{zero} means therefore that, for every pair of open bilateral ideals $\mathfrak{J}\subset A$, $\mathfrak{R}\subset B$ such that $\rho(\mathfrak{J})\subset\mathfrak{R}\subset\mathfrak{R}_0$ and every $(A/\mathfrak{J})$-extension $E$ of $B/\mathfrak{R}$ by $L$, there exist two open ideals $\mathfrak{J}'\subset\mathfrak{J}$, $\mathfrak{R}'\subset\mathfrak{R}$ such that $\rho(\mathfrak{J}')\subset\mathfrak{R}'$ and the inverse image of $E$ by the homomorphism $B/\mathfrak{R}'\to B/\mathfrak{R}$ is \emph{trivial}.

We leave to the reader the analogous definition of inductive limits $\operatorname{Exaltop}_A(B,L)$ from $\operatorname{Exal}_{A/\mathfrak{J}}(B/\mathfrak{R},L)$ and of $\operatorname{Exalcotop}_{A/\mathfrak{J}}(B/\mathfrak{R},L)$ for the case where $A$ is commutative and $B$ an $A$-algebra (resp.\ a commutative topological $A$-algebra).
\end{env}

\begin{env}[18.5.2]
\label{0.18.5.2}
If one has a commutative diagram of continuous homomorphisms of rings
\[
  \xymatrix{
    B'\ar[r] & B\\
    A'\ar[u]\ar[r] & A\ar[u]
  }
\]
one deduces from it canonically two homomorphisms of additive groups
\[
  \operatorname{Exantop}_A(B,L)\to\operatorname{Exantop}_{A'}(B,L)\to\operatorname{Exantop}_{A'}(B',L)
\]
by passage to the inductive limit from \sref{0.18.3.6.1}.

By virtue of the exactness of the $\varinjlim$ functor in the category of commutative groups, the kernel of the homomorphism
\[
  \operatorname{Exantop}_A(B,L)\to\operatorname{Exantop}_{A'}(B,L)
\]
is the inductive limit of the kernels of the homomorphisms
\[
  \operatorname{Exan}_{A/\mathfrak{J}}(B/\mathfrak{R},L)\to\operatorname{Exan}_{A'/\mathfrak{J}'}(B/\mathfrak{R}',L)
\]
where one took for $\mathfrak{J}'$ the inverse image of $\mathfrak{J}$; one denotes this kernel $\operatorname{Exantop}_{A/A'}(B,L)$.
One defines similarly $\operatorname{Exaltop}_{A/A'}(B,L)$ and $\operatorname{Exalcotop}_{A/A'}(B,L)$.
Finally, if one has a homomorphism
\end{env}

\begin{env}[18.5.3]
\label{0.18.5.3}
Given a topological ring $C$ and two topological $C$-bimodules $M$, $N$, one denotes $\operatorname{Hom.\ cont}_C(M,N)$ the additive group of $C$-homomorphisms \emph{continuous} from $M$ into $N$.
\end{env}

\begin{lemma}[18.5.3.1]
\label{0.18.5.3.1}
Let $C$ be a topological ring, $E$, $L$ two topological $C$-bimodules; suppose that the topologies are linear, that $L$ is discrete and annihilated by an open bilateral ideal of $C$.
Then one has a canonical isomorphism
\begin{equation}
\tag{18.5.3.2}
  \label{0.18.5.3.2}
  \varinjlim\operatorname{Hom}_{C/\mathfrak{R}}(E/V,L)\xrightarrow{\sim}\operatorname{Hom.\ cont}_C(E,L)
\end{equation}
where in the first member the inductive limit is taken over the directed set of pairs $(\mathfrak{R},V)$ such that $\mathfrak{R}$ is an open bilateral ideal of $C$, $V$ an open sub-$C$-bimodule of $E$, such that $\mathfrak{R}.L=L.\mathfrak{R}=0$, $\mathfrak{R}.E\subset V$, $E.\mathfrak{R}\subset V$.
\end{lemma}

\begin{proof}
As $C/\mathfrak{R}$ and $E/V$ are discrete, one has canonical homomorphisms $w_{\mathfrak{R},V}:\operatorname{Hom}_{C/\mathfrak{R}}(E/V,L)\to\operatorname{Hom.\ cont}_C(E,L)$ forming an inductive system, hence a homomorphism from the first member of \sref{0.18.5.3.2} to $\operatorname{Hom.\ cont}_C(E,L)$.
As the homomorphism $E/V'\to E/V$ is surjective for $V'\subset V$, it follows from the definition that the homomorphism $\operatorname{Hom}_{C/\mathfrak{R}}(E/V,L)\to\operatorname{Hom}_{C/\mathfrak{R}}(E/V',L)$ for $\mathfrak{R}'\subset\mathfrak{R}$ is injective.
On the other hand, if $u$ is a continuous $C$-homomorphism of $E$ into $L$, its kernel is an open sub-bimodule $V_0$ of $E$, and $\mathfrak{R}_0.E\subset V_0$, $E.\mathfrak{R}_0\subset V_0$, it is clear that $u$ is the canonical image of a $(C/\mathfrak{R}_0)$-homomorphism of $E/V_0$ into $L$, hence \sref{0.18.5.3.2} is surjective.

This being so, the proposition \sref{0.18.3.8} generalizes as follows to topological rings:
\end{proof}

\begin{proposition}[18.5.4]
\label{0.18.5.4}
Let $B$ be a topological ring linearly topologized, $\mathfrak{J}$ a bilateral ideal of $B$, $C=B/\mathfrak{J}$ the quotient topological ring; $\mathfrak{J}/\mathfrak{J}^2$ \textup{(}where $\mathfrak{J}$ is equipped with the topology induced by that of $B$ and $\mathfrak{J}/\mathfrak{J}^2$ with the quotient topology of that of $\mathfrak{J}$\textup{)} is then canonically equipped with a structure of topological $C$-bimodule.
For every discrete $C$-bimodule $L$ annihilated by an open bilateral ideal of $C$, there exists then a canonical isomorphism
\begin{equation}
\tag{18.5.4.1}
  \label{0.18.5.4.1}
  \operatorname{Hom.\ cont}_C(\mathfrak{J}/\mathfrak{J}^2,L)\xrightarrow{\sim}\operatorname{Exantop}_B(C,L).
\end{equation}
\end{proposition}

\begin{proof}
Indeed, for every open $\mathfrak{R}$ of $B$ such that $(\mathfrak{J}+\mathfrak{R})/\mathfrak{J}$ annihilates $L$, one has, by \sref{0.18.3.9}, a canonical isomorphism
\[
  \operatorname{Hom}_{B/(\mathfrak{J}+\mathfrak{R})}((\mathfrak{J}+\mathfrak{R})/(\mathfrak{J}^2+\mathfrak{R}),L)\xrightarrow{\sim}\operatorname{Exan}_{B/\mathfrak{R}}(B/(\mathfrak{J}+\mathfrak{R}),L)
\]
and it suffices to pass to the inductive limit, taking into account \sref{0.18.5.3.1}.
\end{proof}
